% !TEX root = ../main.tex
\chapter{仿真附录：检查清单与流程表}
\label{ch:sim-appendix}

\paragraph{读完本章，你应能}
\begin{itemize}
  \item 理解为什么“收敛、守恒、交叉验证和版本记录”本身就是物理结论的一部分，而不是报告末尾的形式主义附录。
  \item 掌握 RCWA、FDTD 和 FEM 三类数值方法各自最关键的验证逻辑，并知道它们在 Lumerical 与 COMSOL 风格工作流中的典型落点。
  \item 学会把仿真结果分成“探索性、已验证、可发布”三个层级，用明确门槛控制结论升级。
\end{itemize}

\paragraph{读前准备}
建议已掌握 \cref{ch:rcwa,ch:fdtd-fem,ch:workflow,ch:failure-modes} 中关于方法分工、完整工作流和常见失败模式的内容。

\paragraph{阅读主线}
\ChapterModelLevel{L1→L4}
前文已经讲清了 PCSEL 应该怎样建模，但要让模型成为可信证据，还需要最后一层工作，即验证与发布门槛。本附录先说明为什么验证是物理论证本身，而非清单化装饰；随后分别讨论 RCWA、FDTD 和 FEM 的验证逻辑，重点解释每一项检查到底在防什么误判；再把这些方法层面的检查提升到工作流层面，形成从冷腔到器件级的阶段门槛；最后给出一个通用排错决策树，帮助读者在结果异常时按正确顺序排查，而不是一开始就盲目调网格。

\section{为什么验证不是额外工作，而是物理论证本身}
对 PCSEL 这类开放、多尺度、多物理场问题来说，一次仿真从来不只是“给出一个数值”。它同时还暗含了大量前提：边界是否真正开放，网格是否足够细，时间窗是否足够长，参数版本是否一致，模型是否仍处在近似适用范围内。若这些前提没有被验证，那么输出的数值就不能算物理证据，只能算求解器在某组设置下给出的一个候选答案。

这也是为什么本附录不把验证写成孤立 checklist。需要掌握的是每一项检查对应的物理误差来源。例如，能量守恒检查并非为了让表格更完整，而是为了防止把人工边界反射误判成高 $Q$；阵列尺寸收敛是为了确认你研究的确实是有限器件物理，而不是一个偶然选中的尺寸共振点，而非为了“多跑几轮显得更认真”。

\begin{IntuitionBox}
仿真验证的目标是说明当前证据链支持的物理结论在常见数值误差和建模误判下仍保持不变，而非证明求解器能够运行。
\end{IntuitionBox}

\section{通用验证逻辑：先问收敛，再问不变量，再问模型边界}
无论使用哪一种求解器，验证的顺序都不应是随意的。最稳妥的做法是按三层推进。

第一层是\textbf{数值收敛}。也就是检查当离散尺度继续加密时，关键输出是否趋于稳定。这一层防止的是“纯离散误差被误读成物理趋势”。

第二层是\textbf{不变量或守恒量检查}。例如无源周期散射中的功率平衡、模排序对无关归一化的稳健性、以及同一物理结论在不同激励方式下是否保持。这一层防止的是“设置看起来收敛，但其实在解错误问题”。

第三层是\textbf{模型边界检查}。也就是说，即便数值已经收敛，也要问当前模型的物理前提是否仍然成立。比如，二维有效折射率模型对某结构也许已经数值收敛，但若纵向强矢量效应开始主导，它仍然不能支撑最终三维器件结论。

只有这三层都站住，结果才有资格升级。若只通过第一层而没过第二层或第三层，结果通常仍然只能算探索性结果。

\section{RCWA 验证：验证对象是单胞开放散射，而不是某条窄谱线}
RCWA 的优势在于它能高效处理周期单胞的开放散射，但这也意味着它最容易被误读成“只要谱线足够尖就说明模式足够好”。正确的验证逻辑应当围绕四件事展开。

第一，傅里叶截断是否收敛。也就是说，随着保留倒格矢数目增加，关键量例如共振频率、上下出光比例和目标模与竞争模的排序是否稳定。若排序本身仍会随截断阶数变化，那么任何后续结论都应当暂停。

第二，层状切片和几何逼近是否足够准确。对有刻蚀深度、斜侧壁或复杂层栈的结构，RCWA 常需把纵向轮廓分层近似。此时不仅傅里叶阶数要收敛，纵向切片数也要收敛。否则你看到的谱线变化，可能只是几何逼近误差。

第三，功率平衡是否成立。对无源结构，反射、透射与吸收应满足功率守恒。若不守恒，优先应怀疑截断、材料模型或归一化，而不是先解释谱特征。

第四，谱特征是否真和目标模对应得上。一个高质量 RCWA 结果，不应只给一张反射谱，还应能说明某个谱特征与哪类 Gamma 点候选模相对应，以及该特征是否与前面的 PWE/CWT 结论同向。

若把这些检查映射到软件路径上，在 Lumerical 的 RCWA 工作流里，重点通常是傅里叶阶数、入射角/偏振设置和材料层切片一致性；在 COMSOL 的单胞频域 Wave Optics 路径里，虽然实现框架不是 RCWA，但本质上也要检查“单胞周期散射是否随网格、端口和相位周期边界稳定”。求解器名称不同，验证逻辑并没有变。

若某个单胞设计在 RCWA 中显示一条非常窄的反射特征，但将傅里叶阶数从 $N_G$ 提高到 $N_G+2$ 后，竞争模与目标模的排序发生翻转，那么最合理的结论是“当前单胞结论尚未收敛，不能进入下一阶段”，而非“结构具有很强模式选择”。

\section{FDTD 验证：时间窗、PML 和激励源决定你看到的是模还是伪影}
FDTD 的优势在于它直接处理有限阵列、宽带响应和近远场变换，因此在 PCSEL 工作流中尤其适合做有限器件验证。但它也有典型的风险：若时间窗太短、PML 不稳或激励方式不合适，就会把数值伪影误读成模式特征。

第一，网格与时间步必须以关键输出为准，而不是以“默认稳定条件”为准。稳定只意味着算法不会发散，不意味着提取出的远场、$Q$ 或模排序已经可信。必须显式检查近场、远场主瓣和共振位置对网格是否收敛。

第二，PML 不只是一个吸收层参数，而是开放边界定义的一部分。若 PML 距离器件太近、厚度不足或参数与出射角谱不匹配，就会产生人工反射，进而污染衰减率和远场。很多表面上很高的 $Q$，本质上只是边界没吸干净。

第三，激励源决定了你容易看到什么模。若源位置、偏振或频带过于偏向某一类模式，就可能让目标模看起来“天然更强”，但这并不等价于它在自发起振条件下也一定更有优势。因此，在有限阵列 FDTD 中，建议至少用两类激励方式交叉检查，例如脉冲偶极源与更接近目标模对称性的定制源。

第四，若要从时域衰减提取 $Q$，就必须同时报告拟合时间窗和残差。因为 FDTD 输出常包含多个模的叠加，拟合区间选得不同，得到的衰减率可能明显变化。

从软件映射看，FDTD 主要对应 Lumerical FDTD 或 MEEP 一类时域工具。COMSOL 的主力并非 FDTD，而是频域 FEM。因此，如果你的工作流以 COMSOL 为主，就不应强行用“时域直觉”来解释所有频域结果，而应转向本附录后面的 FEM 验证逻辑。

\section{FEM 验证：收敛的不只是网格，还有问题本身的弱形式实现}
频域 FEM 在 PCSEL 建模中很有价值，特别是处理复杂几何、材料色散和多物理场接口时。但 FEM 的验证重点与 FDTD 不同。对 FEM 来说，真正要小心的除了网格大小，还有单元阶次、PML 形式和弱形式实现中的材料定义。

第一，必须区分 $h$-收敛和 $p$-收敛。也就是说，不仅要问网格细化是否稳定，还要问提高单元阶次后结果是否保持一致。对高对比度电磁边界和强曲率几何，区分这两类收敛尤其重要。

第二，开放系统中的本征频率问题常依赖 PML。若 PML 设计不当，本征模可能带入明显伪影或被错误地局域化。因此，PML 厚度、拉伸参数和位置都必须扫一轮，而不是一次设置后就默认可靠。

第三，要尽可能交叉比较本征问题与驱动问题。若某一共振在本征频率研究中存在，那么在对应频率附近的驱动频域响应中也应能看到相容特征。若二者长期对不上，往往说明边界、材料或模识别出了问题。

第四，弱形式中的材料色散和损耗写法必须和正文约定一致。特别是在包含复介电常数、温度依赖材料和各向异性时，只要实现细节稍有不一致，数值就可能“稳定但错误”。

在工程软件上，这一路径最典型的是 COMSOL 的 Wave Optics 频域或本征频率研究。其优势是与 Semiconductor 和 Heat Transfer 共几何、共网格策略更自然；代价则是验证工作不能偷懒，因为复杂耦合越多，弱形式和边界实现的错误就越不容易被肉眼发现。

对 FEM 本征频率问题，出现一个数值上平滑、形状上也像“模式”的场分布，并不自动说明它就是你要的物理模。只要 PML、端口、材料或对称边界设置不当，求解器同样可以给出自洽但无物理意义的伪模。

\section{收敛性检查的目标：不是让数字不变，而是让结论不变}
很多人做收敛分析时，只盯着一个绝对数值，例如共振波长变化不到多少就算通过。这个标准并非无用，但对 PCSEL 来说还不够，因为真正影响设计决策的通常是结论本身是否稳定，而非单个数值。例如：目标模和竞争模的排序是否稳定，远场主瓣是否仍朝同一方向，热回写后工作窗口是否仍然存在。

因此，更成熟的收敛检查应当至少包含两层阈值。一层是数值量收敛阈值，例如频率、$Q$、温升等变化小于某个比例；另一层是结论阈值，即排序、趋势和设计选择不发生翻转。若第一层过了而第二层没过，那么从工程角度看，结果仍然不够稳健，因为它还不能支撑设计决策。

\section{从冷腔到器件级：需要怎样的阶段门槛}
前文 \cref{ch:workflow} 把工作流分成了 L1 到 L4 四层。本附录则把这四层进一步转成“升级门槛”。\cref{tab:stage-gates-appendix} 总结了一个实用版本。

\begin{table}[htbp]
  \centering
  \footnotesize
  \caption{从冷腔到器件级的阶段门槛。只有当前一层证据满足门槛，结果才应升级到下一层结论。}
  \label{tab:stage-gates-appendix}
  \setlength{\tabcolsep}{4pt}
  \begin{tabularx}{\textwidth}{@{}p{1.3cm}p{2.7cm}Y@{}}
    \toprule
    层级 & 主要输出 & 至少应满足的门槛 \\
    \midrule
    L1 结构候选层 & 候选模、对称性、被动损耗趋势 & 单胞或简化模型收敛；功率平衡与边界一致；候选模可与理论图像对上 \\
    L2 光学阈值层 & 阈值增益、有限阵列远场、模排序 & 有限尺寸结果对阵列尺度和开放边界稳定；目标模相对竞争模排序不翻转 \\
    L3 电学映射层 & 载流子分布、J--V、阈值电流形成机制 & 电学边界与参数来源清楚；器件结论不再依赖均匀增益假设 \\
    L4 热-电-光工作层 & 工作窗口、热漂移、模式稳定性 & 温度场回写后结论仍稳健；工作窗口而非单点最优被明确给出 \\
    \bottomrule
  \end{tabularx}
\end{table}

若要把“多物理场已收敛”写得更可复核，还应给出连续两轮闭环迭代的量化差分，例如
\[
\frac{|\Delta I_{\mathrm{th}}|}{I_{\mathrm{th}}}<2\%,
\qquad
|\Delta T_{\max}|<\SI{1}{K},
\qquad
|\Delta \lambda|<\SI{0.02}{nm},
\]
并说明这些门槛对应的是冷腔层、电热层还是完整回写层。否则，“已收敛”仍然只是一个没有门槛的形容词。

这个表最重要的作用，是防止结论越级。只要尚未通过 L3 门槛，就不能输出器件级阈值电流；尚未通过 L4 门槛，就不能输出热稳定工作窗口。这样做看似保守，实际是在节约研究时间，因为它能避免把错误层级的结果带进后续更昂贵的仿真和实验中。

\section{收敛标准不应只写“已收敛”，而应写成门槛}
对真正可复用的 PCSEL 建模记录来说，“结果已收敛”这句话几乎没有信息量。更有用的写法，是先规定当前层最重要的 2--3 个观测量，再规定它们的允许变化范围。例如，单胞 RCWA 层可以把目标模频率、上下辐射比和阈值代理增益作为门槛量；有限阵列全波层可以把目标模频率、$Q$、主瓣角度和主瓣半高全宽作为门槛量；电热光层则可把阈值电流、峰值温升和工作频率漂移作为门槛量。

\begin{table}[htbp]
  \centering
  \footnotesize
  \caption{多物理场链中的最小收敛门槛示例。}
  \label{tab:multiphysics_convergence}
  \begin{tabularx}{\textwidth}{@{}p{2.2cm}p{3.2cm}Y@{}}
    \toprule
    层级 & 建议观测量 & 收敛写法示例 \\ \midrule
    RCWA 单胞层 & 候选模频率、反射谷位置、上下辐射比 & 增加傅里叶阶数后，三者相对变化均小于预设容差 \\
    FDTD/FEM 冷腔层 & 目标模频率、$Q$、远场主瓣角度、主瓣半高全宽 & 网格/PML/阵列尺寸继续加密时，结论不翻转且变化进入平台区 \\
    电学层 & J--V 曲线拐点、载流子峰值位置、串联电阻估计 & 网格和边界再细化时，关键电学量变化不再显著 \\
    热学层 & 峰值温升、热点位置、热阻提取 & 热网格与热边界变化后，热点位置与热阻进入稳定区 \\
    电热光闭环层 & 阈值电流、工作频率漂移、目标模相对竞争模裕量 & 回写再细化后，目标模身份与窗口结论不翻转 \\
    \bottomrule
  \end{tabularx}
\end{table}

\section{网格收敛最好固定成五步流程}
若每次都凭直觉决定“再加密一点”，很容易把收敛检查做成不可复现的个人习惯。更稳妥的最小流程是五步：第一，先选一个当前最关心的可观测量，例如目标模频率、$Q$ 或主瓣角度；第二，只改一个离散参数，如网格尺寸、单元阶次或时间步长；第三，记录两到三个连续加密层级下的变化；第四，确认变化趋势进入平台区，而不是仅偶然变小一次；第五，再切到下一个离散参数。这个顺序的意义，在于避免多个数值旋钮同时变化后无法归因。

\begin{PitfallBox}
多物理场链里最常见的伪收敛，是“光学网格收敛了，但电热回写后结果又翻转”。这是层级升级后可观测量本身变了，而非收敛失败。因此，任何声称“器件级结果已收敛”的说法，都必须写明它是在冷腔层、电学层、热学层还是闭环层收敛。
\end{PitfallBox}

\section{排错决策树：结果异常时按什么顺序查}
当结果异常时，不宜同时修改多个设置并混合归因。更有效的做法是固定一棵决策树。\cref{fig:debug-tree-appendix} 给出一条适合 PCSEL 建模的最小顺序：先查版本和几何，再查边界与单位，再查收敛，最后才讨论是否可能存在真正的新物理。

\begin{figure}[htbp]
  \centering
  \resizebox{0.9\textwidth}{!}{%
  \begin{tikzpicture}[node distance=0.7cm,
    every node/.style={draw, rounded corners, align=center, minimum height=0.85cm, text width=3.6cm}]
    \node (s) {结果异常或跨方法不一致};
    \node (a) [below=of s] {版本、单位、几何是否一致};
    \node (b) [below left=1.0cm and 1.3cm of a, text width=3.3cm] {否\\先修正输入账本};
    \node (c) [below right=1.0cm and 1.3cm of a, text width=3.3cm] {是\\继续检查边界};
    \node (d) [below=of c] {开放边界、端口、完美匹配层、周期条件是否正确};
    \node (e) [below=of d] {网格、时间窗、傅里叶阶数或单元阶次是否收敛};
    \node (f) [below=of e] {不同激励或第二种方法下趋势是否一致};
    \node (g) [below=of f] {若仍异常，再讨论模型边界或新物理};
    \draw[->] (s)--(a);
    \draw[->] (a)--(b);
    \draw[->] (a)--(c);
    \draw[->] (c)--(d);
    \draw[->] (d)--(e);
    \draw[->] (e)--(f);
    \draw[->] (f)--(g);
  \end{tikzpicture}
  }
  \caption[PCSEL 建模排错决策树]{PCSEL 建模的最小排错决策树。}
  \label{fig:debug-tree-appendix}
\end{figure}

这棵树的意义不在于它是否覆盖所有细节，而在于它强制研究者避免一个高频错误：还没确认输入和边界是否一致，就急着把异常结果解释成模式分裂、非线性或热反馈新现象。对开放系统而言，这种过早解释尤其危险，因为数值边界伪影本来就很容易伪装成模态物理。

只保存“最直观的一次结果图”，而不保存失败设置、收敛扫描和版本差异记录，会让团队在几周后完全失去判断力：既不知道当前结果是不是偶然命中，也不知道过去排除过哪些伪影路径。对 PCSEL 这类多阶段工作流，这种信息丢失的代价极高。

\section{怎样把验证结果写成可发布报告}
当所有验证都完成后，最好的报告方式并非把几十张收敛图全部塞进正文，而是把它们整理成三层证据：

第一层是“当前结论依赖什么模型和近似”。这让读者知道结论的物理边界。

第二层是“当前结论通过了哪些数值门槛”。这让读者知道结果不只是理想图，而是经得起反证的结果。

第三层是“当前结论仍受哪些不确定因素限制”。这让读者知道结果是否已达到可发布级别，还是仍然停留在探索性级别。

对软件路径而言，也应采取同样写法。例如，不必写成“在某软件里点了哪些选项”，而应写成“在周期单胞开放散射模型中已验证傅里叶/网格收敛，在有限阵列开放全波模型中已验证 PML 与阵列尺寸收敛，在器件模型中已验证电热回写后结论不翻转”。这样写既保持与求解器实现解耦，又保留了足够的技术可追溯性。
\section{数值不确定度量化与报告规范}
基于 \cref{ch:method-compare} 中的误差分析理论，本节给出可直接用于项目实践的数值不确定度量化流程和报告模板。

\subsection{各方法的主要误差源}
\label{sec:numerical-error-per-method}
多方法互证的另一个关键维度是理解并量化每种方法的数值误差来源。只有清楚地知道每个数字背后的不确定度，才能做出可靠的工程判断。本节系统梳理各类方法的主要误差源、收敛行为和交叉验证策略。


\paragraph{先统一本章收敛门槛口径（探索级/本书算例/发布级）}
为避免把“调参初筛阈值”误读成“发布级硬标准”，本章后续阈值统一按三层门槛解释：探索级用于快速筛查方向；\textbf{[本书算例]}用于本教材 benchmark 对齐；发布级用于论文或项目验收。若项目硬件、几何或误差预算变化，应重新定义门槛，而不是机械照搬。

\begin{table}[htbp]
  \centering
  \footnotesize
  \caption{本章收敛门槛的三层口径（方法学文献 + 本书 benchmark 约定）。}
  \label{tab:method_threshold_tiers}
  \setlength{\tabcolsep}{4pt}
  \begin{tabularx}{\textwidth}{@{}p{3.0cm}p{3.1cm}p{3.3cm}Y@{}}
    \toprule
    指标 & 探索级（快速筛查） & 本书算例（章节默认） & 发布级（建议） \\
    \midrule
    PWEM 频率收敛 & 相邻 $N_G$ 频率漂移 $<1\%$ & 相邻 $N_G$ 频率漂移 $<0.5\%$ & 相邻 $N_G$ 频率漂移 $<0.2\%$ 并给收敛曲线 \\
    RCWA 共振峰漂移 & 阶数扫描后峰位变化 $<2\,$nm & 峰位变化 $<0.5\,$nm & 峰位变化 $<0.2\,$nm 且功率守恒误差 $<10^{-3}$ \\
    FDTD/FEM 共振频率一致性 & 同数量级一致 & 相对偏差 $<1\%$ & 相对偏差 $<0.5\%$ 并给网格与时间窗敏感性 \\
    模式排序一致性 & 目标模不翻转 & 目标模不翻转 + 竞争模次序稳定 & 目标模不翻转并给最坏样本与误差条 \\
    \bottomrule
  \end{tabularx}
\end{table}

\subsubsection{PWEM 的误差来源与控制}
PWEM 作为本征模求解器，其主要误差来自三个方面：

\paragraph{平面波截断误差}
PWEM 将介电常数分布展开为傅里叶级数，截断阶数$N_G$直接影响精度。典型的收敛行为是：随着$N_G$增加，本征频率按幂律收敛。对高对比度结构（如空气孔/半导体界面），需要更高的$N_G$才能达到与低对比度结构相同的精度。

\begin{NumericalErrorBox}
\textbf{PWEM 收敛性检查清单}
\begin{itemize}
  \item \textbf{频率收敛}：目标模式的归一化频率$\omega a/(2\pi c)$随$N_G$的变化应小于$10^{-4}$
  \item \textbf{带隙位置收敛}：带隙边缘的频率漂移应小于带隙宽度的 1\%
  \item \textbf{场分布收敛}：连续两次$N_G$增加的场均方根差异应小于 1\%
  \item \textbf{对称性保持}：简并模的频率分裂应远小于它们与邻近模式的间距
\end{itemize}
典型经验值：对空气桥型 PC，$N_G \geq 11\times 11$；对衬底支撑型 PC，$N_G \geq 15\times 15$。
\end{NumericalErrorBox}

\paragraph{k 空间采样误差}
对色散曲线的绘制，$\kvec$空间的采样密度会影响带结构的平滑度和极值点的准确性。建议在布里渊区高对称点附近加密采样，以确保带边位置的精度优于$0.5\%$。

\paragraph{有效折射率近似误差}
二维 PWEM 通常用有效折射率$n_{\text{eff}}$代替真实的三维波导结构。这会引入系统性误差，特别是对垂直损耗和模式体积的计算。定量估计时，应将二维 PWEM 结果与三维 RCWA 或 FEM 进行基准比对，典型偏差范围在$5\%-15\%$之间。

\subsubsection{CWT 的模型降阶误差}
CWT 通过将无穷维 Maxwell 方程投影到少数主导波上来实现模型降阶。其误差主要来自 truncation 和耦合系数近似。

\paragraph{主导波选择的完备性}
CWT 的精度取决于所选基底是否能充分描述真实物理。一个实用的检验方法是逐步增加波的数目，观察关键输出（如阈值增益、辐射效率）的相对变化：
\begin{equation}
\varepsilon_N = \left|\frac{X_{N+1}-X_N}{X_N}\right|,
\label{eq:cwt_convergence_metric}
\end{equation}
其中$X_N$是用$N$个波计算的目标量。当$\varepsilon_N < 1\%$时可认为收敛。

\paragraph{耦合系数的计算精度}
CWT 中的耦合系数$\kappa$通常通过对单胞积分获得：
\begin{equation}
\kappa = \frac{\omega^2}{2V}\int_V \Delta\varepsilon(\rvec)\,\E_i^*(\rvec)\cdot\E_j(\rvec)\,\dd^3 r.
\label{eq:cwt_coupling_integral}
\end{equation}
数值积分时的网格密度和几何逼近会直接影响$\kappa$的精度。建议做法是：先用高精度全波方法（RCWA/FEM）计算单胞场，再用这些场校准 CWT 的耦合系数。

\begin{CrossValidationBox}
\textbf{CWT 与 RCWA 的交叉验证协议}
\begin{enumerate}
  \item 用相同单胞几何分别运行 CWT 和 RCWA
  \item 比较两者的共振频率偏差：$|\lambda_{\text{CWT}}-\lambda_{\text{RCWA}}|/\lambda_{\text{RCWA}} < 2\%$
  \item 比较辐射效率的趋势一致性（不必逐点相同）
  \item 比较参数扫描的灵敏度符号（如对孔径变化的响应方向）
\end{enumerate}
若以上任一条件不满足，应返回检查 CWT 的主导波选择或 RCWA 的傅里叶阶数。
\end{CrossValidationBox}

\subsubsection{RCWA 的傅里叶收敛与 Gibbs 现象}
RCWA 的核心是将电磁场和介电函数展开为 Floquet-Bloch 级数。其主要误差源包括：

\paragraph{傅里叶阶数不足导致的 Gibbs 振荡}
对具有高对比度界面的结构，介电函数的傅里叶级数在不连续处会出现 Gibbs 振荡。这会导致：
\begin{itemize}
  \item 共振峰位置的系统性偏移（通常向长波方向）
  \item 虚假的旁瓣或共振分裂
  \item 功率守恒 violation（反射 + 透射$\neq 1$）
\end{itemize}

\begin{NumericalErrorBox}
\textbf{RCWA 收敛性检查清单}
\begin{itemize}
  \item \textbf{傅里叶阶数扫描}：从$N_G=10$开始，每次增加 5，直到共振波长变化$<0.5\,$nm
  \item \textbf{功率守恒}：$|R+T+A-1| < 10^{-3}$（无源结构）
  \item \textbf{ reciprocity 检查}：正向和反向入射应给出相同的透射谱
  \item \textbf{层厚收敛}：对有刻蚀深度的结构，纵向切片数应使得每片厚度$<\lambda/(20 n)$
\end{itemize}
对深孔结构（纵横比$>5:1$），可能需要$N_G > 30$才能收敛。
\end{NumericalErrorBox}

\paragraph{坐标变换技术的必要性}
对圆柱形空气孔等曲面边界，标准的阶梯近似会引入显著的几何误差。此时应采用因子化规则（factorization rules）或曲线坐标系 RCWA 来提高精度。一个简单的诊断是：若增加$N_G$后共振波长仍在单调漂移（而非震荡收敛），则很可能需要改进几何表述。

\subsubsection{FDTD 的时间域误差累积}
FDTD 通过在时间和空间上离散 Yee 网格来演化 Maxwell 方程。其误差特性与其他方法显著不同。

\paragraph{数值色散误差}
Yee 网格中的波传播速度依赖于波长与网格尺寸之比。这会导致：
\begin{equation}
v_{\text{num}} = c\left[1 - \frac{1}{24}(k\Delta x)^2 + O((k\Delta x)^4)\right],
\label{eq:fdfd_numerical_dispersion}
\end{equation}
其中$k=2\pi/\lambda$。为将相速度误差控制在 1\%以内，需要$\Delta x < \lambda/(10n)$。对 PCSEL 这种多波长干涉的结构，建议使用更保守的标准：$\Delta x < \lambda/(15n)$。

\paragraph{时间窗长度与 Q 值提取误差}
从高$Q$谐振器的衰减曲线提取$Q$值时，时间窗长度$T_{\text{win}}$必须满足：
\begin{equation}
T_{\text{win}} > \frac{10 Q}{\omega_0},
\label{eq:fdfd_time_window_requirement}
\end{equation}
否则拟合得到的衰减速率会有显著偏差。一个常见的症状是：延长仿真时间后，提取的$Q$值持续增大而不饱和。

\paragraph{PML 反射引起的伪影}
完美匹配层（PML）的理论反射率为零，但数值实现中会有微小反射。对超高$Q$模式（$Q>10^6$），即使是$10^{-6}$量级的 PML 反射也会污染结果。缓解策略包括：
\begin{itemize}
  \item 增加 PML 层数（典型值：12-16 层）
  \item 使用渐变电导率剖面（polynomial grading）
  \item 将 PML 放置在远离强场区域的位置（至少$\lambda/2$）
  \item 若问题本身是无限周期单胞，可在横向采用 periodic boundary + Bloch 条件；若是有限器件，应保留开放边界（PML），并通过增大计算域、提高 PML 厚度/阶次与参数扫描抑制伪反射
\end{itemize}

\begin{CrossValidationBox}
\textbf{FDTD 与 FEM 的交叉验证协议}
\begin{enumerate}
  \item 在同一有限尺寸器件上分别运行 FDTD（时域）和 FEM（频域）
  \item 比较共振频率：偏差应$<1\%$
  \item 比较$Q$值：在对数尺度上一致即可（如两者都在$10^4$量级）
  \item 比较近场图案的重叠积分：$|\langle E_{\text{FDTD}}|E_{\text{FEM}}\rangle|^2 > 0.9$
  \item 比较远场主瓣方向和角宽：峰值角度偏差$<2^\circ$
\end{enumerate}
注意：由于 FDTD 是宽带激励而 FEM 是单频求解，前者可能激发更多杂散模，因此近场重叠略低是可接受的。
\end{CrossValidationBox}

\subsubsection{FEM 的网格自适应与弱形式误差}
FEM 通过变分原理将微分方程转化为线性代数问题。其误差估计最为成熟。

\paragraph{h- refinement 与 p-refinement}
FEM 提供两种收敛途径：
\begin{itemize}
  \item \textbf{h-精化}：减小网格尺寸$\Delta h$，固定单元阶次$p$
  \item \textbf{p-精化}：提高单元多项式阶次$p$，固定网格
\end{itemize}
对光滑解，p-精化的指数收敛速度远超 h-精化。但对具有奇点（如尖锐金属边角）的问题，h-精化更为稳健。

\paragraph{残差型后验误差估计}
现代 FEM 软件通常提供基于残差的误差指示器：
\begin{equation}
\eta_K = \|R(u_h)\|_{L^2(K)} + \|J(u_h)\|_{L^2(\partial K)},
\label{eq:fem_residual_estimator}
\end{equation}
其中$R(u_h)$是单元内部残差，$J(u_h)$是通量跳跃。 adaptive mesh refinement (AMR) 会根据$\eta_K$自动加密高误差区域。

\begin{NumericalErrorBox}
\textbf{FEM 收敛性检查清单}
\begin{itemize}
  \item \textbf{自由度扫描}：绘制目标量 vs 自由度数，确认已进入渐近平台区
  \item \textbf{单元阶次提升}：从二次单元升到三次单元，目标量变化应$<1\%$
  \item \textbf{PML 敏感性}：移动 PML 边界或改变 PML 参数，本征频率漂移应$<0.1\%$
  \item \textbf{能量范数收敛}：$\|u-u_h\|_E/\|u\|_E$应随自由度增加而单调下降
\end{itemize}
对本征频率问题，还需额外检查：虚部（损耗）的收敛速度通常慢于实部（频率）。
\end{NumericalErrorBox}


\subsection{最小必要收敛扫描组合}
对 PCSEL 多方法工作流，不建议对所有参数都做全面扫描（那样成本过高），而是聚焦于对最终结论最敏感的三个维度：

\begin{table}[htbp]
  \centering
  \footnotesize
  \caption{各方法的最小必要收敛扫描组合。}
  \label{tab:min-convergence-scan}
  \begin{tabularx}{\textwidth}{@{}p{2cm}p{3.5cm}X@{}}
    \toprule
    方法 & 必扫参数 & 通过标准 \\
    \midrule
    PWEM & 平面波截断数 $N_G$ & 目标模频率变化 $< 0.5\%$，简并模分裂 $<$ 邻近模间距的 10\% \\
    RCWA & 傅里叶阶数 + 纵向切片数 & 共振波长漂移 $< 0.5\,$nm，功率不平衡 $< 10^{-3}$ \\
    FDTD & 网格密度 + 时间窗长度 & $Q$值变化 $< 5\%$，远场主瓣角度漂移 $< 2^\circ$ \\
    FEM & 网格自由度 + 单元阶次 & 本征频率实部变化 $< 0.1\%$，虚部变化 $< 10\%$ \\
    CWT & 主导波数量 & 阈值增益相对变化 $< 2\%$ \\
    \bottomrule
  \end{tabularx}
\end{table}

实际操作时，建议采用"两步粗扫 + 一步精扫"策略：先在较大步长下确定大致平台区，然后在平台区内做一次精细扫描，最后用 Richardson 外推或直接取平台区的平均值作为最佳估计。这样可以避免陷入局部平台或误判收敛趋势。

\subsection{不确定度的合成与传递}
对多步骤工作流，前一级仿真输出的不确定度会成为下一级输入的误差源。设第$i$级输出的相对不确定度为$\delta_i$，则最终结果的总相对不确定度可用平方和根 (RSS) 估算：
\begin{equation}
\delta_{\text{total}} = \sqrt{\sum_{i=1}^N w_i^2 \delta_i^2},
\label{eq:uncertainty rss}
\end{equation}
其中$w_i$是第$i$级对最终结果的灵敏度权重。对缺乏严格灵敏度分析的情况，保守做法是取$w_i=1$。

\begin{NumericalErrorBox}
\textbf{典型 PCSEL 工作流的不确定度预算示例}
假设一个完整的冷腔→阈值电流预测链条包含以下步骤：
\begin{itemize}
  \item PWEM 带结构：$\delta_1 \approx 0.5\%$（频率定位）
  \item CWT 耦合系数：$\delta_2 \approx 3\%$（辐射损耗提取）
  \item 速率方程阈值：$\delta_3 \approx 5\%$（载流子输运模型）
  \item 热回写迭代：$\delta_4 \approx 2\%$（温度依赖折射率）
\end{itemize}
按 RSS 合成，总不确定度约为$\delta_{\text{total}} \approx \sqrt{0.5^2+3^2+5^2+2^2}\% \approx 6.2\%$。这里应明确该结果属于\textbf{[本书算例]}，不应直接外推到任意器件。它提醒我们：即使每一步都已“收敛”，最终阈值电流预测仍可能存在数个百分点级别的不确定度。因此，声称“仿真阈值为 150 mA”不如写成“仿真阈值为 $150 \pm 10\,$mA”更符合报告规范。
\end{NumericalErrorBox}

\begin{WhatCouldGoWrong}
一个常见但危险的误区是：在前几步用较粗糙的网格快速筛选，期望在最后一步用精细网格"纠正"所有误差。这种做法忽略了误差的非线性放大效应——早期阶段的定性误判（如模式排序错误）无法通过后期的高精度计算挽回。正确的策略是：在每一层都确保数值误差不至于改变定性结论。
\end{WhatCouldGoWrong}

\subsection{交叉验证的可接受偏差范围}
当使用两种独立方法进行交叉验证时（如 FDTD vs FEM），需要事先定义什么是"一致"。表~\ref{tab:xval-tolerance}给出的阈值是\textbf{[本书建议门槛]}，用于教学和工程筛查，不代表社区统一标准。

\begin{table}[htbp]
  \centering
  \footnotesize
  \caption{跨方法交叉验证的可接受偏差范围（\textbf{[本书建议门槛]}，可按项目重定义）。}
  \label{tab:xval-tolerance}
  \setlength{\tabcolsep}{3pt}
  \begin{tabularx}{\textwidth}{@{}p{2.8cm}p{3.2cm}X@{}}
    \toprule
    待比较量 & 宽松容差 (探索级) & 严格容差 (发布级) \\
    \midrule
    共振波长/频率 & $< 5\,$nm 或 $< 3\%$ & $< 1\,$nm 或 $< 0.5\%$ \\
    $Q$值 & 同数量级即可（\textbf{[探索级]}） & 相对偏差 $< 20\%$（\textbf{[发布级]}） \\
    远场主瓣角度 & $< 5^\circ$ & $< 2^\circ$ \\
    近场重叠积分 & $> 0.8$ & $> 0.9$ \\
    阈值增益 & $< 10\%$ & $< 3\%$ \\
    模式排序 & 完全一致 & 完全一致（硬性要求） \\
    \bottomrule
  \end{tabularx}
\end{table}

这些门槛的来源可追溯到 RCWA/FDTD/FEM 的经典收敛文献与有限尺寸敏感性讨论\cite{li1996,moharam1995stablercwa,wang2025finitesizepcsel}；但具体项目应按器件波段、几何复杂度和实验误差预算重定义后再用于发布级结论。

\begin{WhatCouldGoWrong}
一个常见误区是用“宽松容差”级别的吻合去支撑“发布级”强结论。例如，若两种方法的 $Q$ 值只在同一数量级（如 $10^4$ vs $2\times10^4$），就只能归类为\textbf{[工程经验]/探索级}一致，而不能宣称\textbf{[文献值]/发布级}定量可靠。正确做法是先声明采用的容差等级，再限定结论强度。
\end{WhatCouldGoWrong}

\subsection{可重复性报告的最小信息集}
为确保他人能够复现你的数值结果，建议在论文或技术报告的附录中包含以下信息：

\begin{ReportingStandardBox}
\textbf{PCSEL 数值仿真的最小可重复信息清单}
\begin{enumerate}
  \item \textbf{几何参数}：晶格常数、孔径/占空比、刻蚀深度、层厚（附公差或加工变异范围更佳）
  \item \textbf{材料模型}：折射率数据来源（实测/文献/拟合）、色散公式类型（Cauchy/Sellmeier/Drude-Lorentz）、是否包含虚部
  \item \textbf{离散设置}：网格类型（结构化/非结构化）、关键区域网格尺寸、单元阶次（FEM）、时间步长（FDTD）
  \item \textbf{边界条件}：PBC/PML/PEC 的位置和参数、PML 层数和吸收剖面
  \item \textbf{收敛证据}：至少提供一组关键量的收敛曲线（可在 Supplemental Material 中）
  \item \textbf{软件信息}：求解器名称、版本号、关键插件或自定义脚本
  \item \textbf{硬件资源}：CPU/GPU型号、内存、单次仿真耗时（有助于他人评估复现成本）
\end{enumerate}
示例表述（可放入 Methods 部分）："RCWA 仿真使用 S4 求解器 (v1.8)，傅里叶阶数$N_G=21$，纵向切片数 12 层。网格收敛测试见 Supplementary Fig. SX，共振波长标准偏差为$0.6\,$nm。"
\end{ReportingStandardBox}

若期刊不允许长篇方法细节，可将完整收敛数据和脚本上传至公开代码仓库（如 GitHub、Zenodo 或 figshare），并在正文中提供 DOI 链接。这已成为计算光子学领域的最佳实践，也是许多高水平期刊的强制要求。
\section*{回看本章}
仿真验证是物理论证的一部分，而非报告末尾的形式性清单。RCWA 需要证明单胞开放散射的截断和功率平衡可靠，FDTD 需要证明时间窗、PML 和激励源没有把伪影当成模，FEM 需要证明网格、单元阶次和弱形式实现不会制造伪模。更高一层地说，所有这些方法验证最终都要服务于工作流门槛：只有当 L1 到 L4 的证据依次满足基本约束，结论才应升级。可靠的 PCSEL 建模，不只是“求解器完成运行”，而是“结论有门槛、异常有决策树、结果有可发布级别”。

\section*{练习题}
\begin{enumerate}
  \item \basicq 为一个你当前正在研究的 PCSEL 结构，分别写出 RCWA、FDTD 或 FEM 中最关键的三项验证门槛，并说明每一项各自防止什么误判。

  \item \basicq 结合 \cref{tab:stage-gates-appendix}，说明为什么一个已经通过 L2 的结构仍然不能直接宣称器件级阈值电流优越。

  \item \midq 试着把 \cref{fig:debug-tree-appendix} 改写成适合你自己项目的排错流程，并加入你最常遇到的一类异常症状。

  \item \hardq 设计一份一页纸验证报告模板，使其同时适用于以 Lumerical 为主和以 COMSOL 为主的工作流。
\end{enumerate}

\section*{延伸阅读}
\begin{itemize}
  \item 建议与 \cref{ch:workflow,ch:failure-modes} 对照阅读，把本附录直接转化成项目级发布门槛和调试清单。 阅读时应同时核对本章变量在后续模型中的定义和适用范围。

  \item 建议与 \cref{ch:worked-example,ch:failure-modes} 对照阅读，把本附录中的验证门槛、例题中的执行步骤和失败模式中的反向诊断锁成一条完整证据链。 阅读时应同时核对本章变量在后续模型中的定义和适用范围。
\end{itemize}

